\chapter[Documentación Técnica \textit{\textbf{Input Collector}}]{Módulo \textit{\textbf{Input Collector}} - Documentación Técnica}
\label{chap:anexo_input_collector}

\section{Introducción al Módulo}
\label{sec:intro_input_collector}

El módulo \textit{\textbf{Input Collector}} proporciona capacidades especializadas para la captura, procesamiento y visualización de eventos de entrada desde múltiples dispositivos en tiempo real. Su desarrollo se orientó específicamente hacia la resolución de los desafíos asociados con la sincronización temporal precisa y el manejo de dispositivos heterogéneos.

\subsection{Objetivos Técnicos}
\label{subsec:objetivos_input_collector}

El desarrollo del módulo se centró en objetivos técnicos específicos:

\begin{itemize}
    \item \textbf{Captura Universal}: Arquitectura extensible que unifica el manejo de dispositivos diversos
    \item \textbf{Precisión Temporal}: Captura y procesamiento con latencia mínima y marcado temporal preciso
    \item \textbf{Accesibilidad Múltiple}: Interfaces diversas que facilitan diferentes patrones de uso
    \item \textbf{Visualización Dinámica}: Representaciones gráficas adaptativas según el dispositivo activo
\end{itemize}

\section{Documentación Técnica Detallada}
\label{sec:doc_tecnica_input_collector}

La documentación técnica completa del módulo, incluyendo especificaciones de arquitectura, patrones de diseño, interfaces de programación, casos de uso detallados y guías de implementación, se encuentra disponible en el siguiente documento:

\includepdf[pages=-]{imaxes/input_collector_doc.pdf}

\section{Casos de Uso Principales}
\label{sec:casos_uso_input}

\subsection{Generación de \textit{\glspl{dataset}} \acrshort{bci}}
\label{subsec:generacion_datasets_bci}

El caso de uso principal del módulo es la captura de eventos de entrada con \textit{\glspl{timestamp}} precisos que faciliten la sincronización posterior con señales cerebrales. Esta capacidad permite generar \textit{\glspl{dataset}} de alta calidad para entrenar sistemas \acrshort{bci}, donde la correlación temporal entre intenciones (detectadas mediante \acrshort{eeg}) y acciones físicas del usuario es fundamental para el desarrollo de modelos predictivos efectivos.

\subsection{Análisis de Comportamiento de Usuario}
\label{subsec:analisis_comportamiento_usuario}

El módulo facilita el análisis detallado de comportamientos de usuario durante tareas específicas, como partidas de videojuegos o interacciones con interfaces complejas. Los datos capturados permiten identificar patrones de uso, evaluar estrategias de interacción, y correlacionar comportamientos observados con estados mentales detectados mediante técnicas neurofisiológicas.

\subsection{Validación y \textit{Testing} de Controladores}
\label{subsec:validacion_testing}

Las capacidades de visualización en tiempo real del estado completo del dispositivo, análisis de latencia de respuesta con marcado temporal preciso, y calibración de componentes analógicos proporcionan herramientas especializadas para validación técnica de dispositivos de entrada.

\subsection{Investigación en Experiencia de Usuario}
\label{subsec:investigacion_experiencia}

El seguimiento detallado de interacciones facilita estudios de usabilidad mediante análisis de patrones de movimiento, generación de datos para análisis de comportamiento de usuario, y evaluación de accesibilidad con monitorización de patrones de uso.